Rebuttal (Reviewer ByWV)

We thank the reviewer for the careful reading and constructive feedback. We respond below in the same order as the review.

Q1 (Additional experiments)

We agree that broader empirical coverage would strengthen the paper. Subject to the page budget, we will add further experiments where feasible, emphasize variance over random seeds (and, where relevant, over frozen-feature draws) to quantify robustness, and include a direct comparison in which a full-rank baseline is trained from i.i.d.\ initialization under the same optimizer and schedule as the RF-LR models. This is the fairest protocol for testing whether avoiding collapse translates into consistent empirical gains relative to a full-rank reference.

Q2 (Convergence guarantee under i.i.d.\ initialization)

Our main theoretical statement is conditional: if the mean-field dynamics converges to a limit, that limit achieves global optimality within the RF-LR parametrized class under the stated assumptions. We do not claim a general unconditional convergence rate or finite-time guarantee for the training dynamics themselves. We will foreground this distinction in the camera-ready and avoid wording that could be read as an unconditional convergence theorem.

Q3 (What is new relative to Nguyen \& Pham, 2023)

We will add a dedicated "What is new" paragraph and a short roadmap figure in prose. Conceptually, we study low-rank structure together with random frozen features; Nguyen \& Pham (2023) supply the general mean-field analytical vocabulary in which we embed that architecture. To our knowledge, this is the first work that carries their multilayer mean-field program into the RF-LR setting (frozen mixing, $r$ channel bundles, and the induced coupling estimates) rather than reproducing their original full-rank i.i.d.\ configuration. The manuscript is self-contained on purpose: after the background required to parse Nguyen \& Pham (2023), we write the full argument so that every object is defined before convergence claims are made. The opening part of the appendix is devoted to existence, uniqueness, and well-posedness of the mean-field objects, because global-convergence statements are only meaningful for dynamics that are proved to exist. In our view, that well-posedness block is largely standard once the RF-LR scaling is fixed; the substantive new difficulty is the global-optimality step after low-rank channel aggregation enters the backward Lipschitz chain, namely replacing the full-rank contraction route with the $\|W^{(\ell)}\|_{\infty,1}$-controlled channelwise estimates and the associated weighted coupled-$L^1$ assumptions. That step is genuinely specific to RF-LR and is not a line-by-line repetition of Nguyen \& Pham (2023).

More broadly, one motivation for this work is to understand how low-rank parametrizations reshape training and the loss landscape relative to full-rank training across scaling regimes (mean-field, feature learning, kernel-like limits). This paper is one systematic mean-field route to that question in the RF-LR architecture. We will make that scope clearer without overstating what the current theorems imply beyond this regime.

Q4 (Assumptions in main text vs.\ appendix)

We agree. We will move a compact assumption block into the main text and add a short discussion for each item: what it buys in the argument, and whether it is structural (architectural richness) or technical (regularity, couplings).

Q5 (Typos and notation)

We agree and will correct the reported issues in the camera-ready, including the extra space after the question mark, the broken fragment around "substantially," and the vector notation typo.

Q6 (Global minimizer versus low-rank feasible set)

We agree with the reviewer’s logic. Low-rank and full-rank parametrizations generally induce different feasible sets in function space, so one must not equate global optimality in one problem with global optimality in the other. Our statements concern global optimality within the RF-LR class, conditional on mean-field convergence. Empirical curves should therefore be read as evidence about training behavior and landscape effects under that architecture, not as a claim that every full-rank global minimizer is matched or attainable under rank constraints. We will make this interpretive framing clearer in the revision so that the figures are not over-read as full-rank global-optimality claims.

Q7 (Momentum selection across rank in the figures)

Momentum and learning-rate values in the paper were obtained from optimizer sweeps over a prescribed grid; we will report the exact ranges and selection rule in the camera-ready. In the main plots we will compare ranks only under matched momentum and matched auxiliary hyperparameters so that rank is not confounded with per-rank tuning. Exploratory momentum-versus-rank comparisons, when shown, will be moved to appendix or supplementary material.

Q8 (Relation to Zhang et al., 2023)

We agree that, at the architectural level, the model is essentially equivalent to the multi-component template of Zhang et al.\ (2023). Our contribution is at the mean-field parametrization and optimization-analysis level: we derive the RF-LR mean-field ODE system for the $r$ channel bundles, prove a conditional global-optimality result when those dynamics converge, and formalize the channel-specialization mechanism. Zhang et al.\ do not provide this mean-field training analysis.

Q9 (Nguyen \& Pham (2023) convergence assumptions)

We will add a concise boxed explanation in the manuscript. In brief, Nguyen \& Pham (2023) impose convergence to a limit through weighted coupled-$L^1$ functionals under couplings $\pi_t$ of finite-width laws with their mean-field limits. These are not statements of bare moment convergence: the integrands pair a nonnegative weight built from downstream magnitudes along the backpropagation chain with a coupled gap between finite-$t$ and limit trajectories, so the condition is tailored to the Lipschitz constants that appear in their global-optimality argument. They also require time-regularity (essential-supremum decay of selected weight velocities) to pass finite-$t$ gradient-flow identities to the limit; this part is proof-relevant, not cosmetic. Heuristically, prefactors such as $|\bar w_2(C_1)|$ on the $w_1$-gap reflect sensitivity: large downstream mass makes $w_1$-errors more consequential for the output, so convergence is measured in a topology aligned with the backward map rather than in unweighted $L^2$ of raw weights.

Our appendix writes the RF-LR analogue as weighted integral conditions such as \texttt{eq:conv-w1}--\texttt{eq:conv-w2}, with the low-rank modification that channel aggregation enters through $\max$- or $\sum$-type control over the $r$ first-layer channels. Thus the same weighted coupled-$L^1$ template is retained, but adapted to $r$ bundles.

For clarity in the camera-ready, we will summarize the template as follows. In the two-layer case, one asks for a coupling $\pi_t$ such that
$$
\mathbb{E}_{\pi_t}\Bigl[\bigl|\bar w_2(C_1)\bigr|\,\bigl|w_1(t,C_1')-\bar w_1(C_1)\bigr|+\bigl|w_2(t,C_1',1)-\bar w_2(C_1)\bigr|\Bigr]\to 0
$$
possibly together with a condition such as $\operatorname*{ess\,sup}_t |\partial_t w_2(t,C_1,1)|\to 0$ to commute limits along the flow. In multilayer settings, the same pattern appears with downstream products of limiting weights multiplying the coupled gaps. In RF-LR, the first-layer gap is controlled channelwise and then aggregated across the $r$ channels.

We will also note that Wasserstein-$4$ convergence is a sufficient technical route, not the definition of the assumption: with suitable moment bounds, H\"older yields the weighted coupled-$L^1$ quantities from a $W_4$ control on the coupled particle system.

Q10 (SGD versus mean-field ODE)

Yes. Following Nguyen \& Pham (2023), the mean-field ODE is the rigorous continuous-time scaling limit of SGD when the step size is sent to zero together with the mean-field scaling; see in particular their Proposition~23 in the discrete-to-continuous discussion. Our quantitative finite-width statement tracks both the $\mathcal{O}(1/\sqrt{n})$ sampling error and the $\mathcal{O}(\sqrt{\varepsilon})$ discretization error alongside the usual Gr\"onwall factors; we will make this alignment with their SGD-to-flow methodology more explicit in the revision.
